1\section{Introduction}

Martin-L\"of's inductive definition system \cite{Martin-Lof-1971} is
widely used for formalizing inductive definitions.
Cyclic proof systems \cite{Brotherston06,Brotherston11a} are another proof system to formalize
inductive definitions and have been actively 
studied \cite{das20,das21,das22,das23a,das23,das26,berardi17,berardi19,brotherston07,brotherston08,brotherston11b,brotherston12,brotherston17b,brotherston17a,brotherston17,oda25,brotherston25},
since cyclic proof systems may work in a better way for proof search.

To prove soundness of cyclic proof systems or
to reduce a cyclic proof to an ordinary inductive proof,
a proof transformation from cyclic proofs to inductive proofs is 
important \cite{Berardi17a,Berardi18,berardi19}.

The first contribution of this paper gives intuitionistic ranking function synthesis
for fan-in-free case in intuitionistic meta-theory.
The second contribution is that
by using it
this paper prove intuitionistic proof transformation
from intuitionistic first-order cyclic proof system with Heyting arithmetic
into intuitionistic first-order inductive proof system with Heyting arithmetic.

The classical first-order cyclic proof system $\CLKIDomega$ and 
the classical first-order inductive proof system $\LKID$ are defined 
in \cite{Brotherston11a}.
As defined in \cite{Berardi18},
we define systems $\CLJIDomega$ and $\LJID$ 
as intuitionistic versions of $\CLKIDomega$ and $\LKID$.
As defined in \cite{Berardi18},
we also define systems $\CLJIDHA$ and $\LJIDHA$ as
the system $\CLJIDomega$ with Heyting arithmetic $\HA$
and the system $\LJID$ with Heyting arithmetic $\HA$.

First, we compare our results for proof transformation
with existing work \cite{Berardi18} and
\cite{Ikebuchi25} closely related to this paper.

% (1)
In \cite{Berardi18},
they extended
the proof transformation in \cite{Berardi17}
to intuitionistic systems $\CLJIDHA$ and $\LJIDHA$,
by replacing the proof of Podelski-Rybalchenko theorem in  \cite{Berardi17}
by an intuitionistic Podelski-Rybalchenko theorem.

In this paper,
we will give
another proof transformation from
the intuitionistic system $\CLJIDHA$ and 
to the intuitionistic system $\LJIDHA$.
The result is similar to \cite{Berardi18},
but in this paper,
we use a ranking function
to prove the correctness of the proof transformation
and on the other hand
the paper \cite{Berardi18} used 
Kleene-Brouwer theorem
to prove the correctness of the proof transformation.
The meta-theory used for proofs in both papers is
intuitionistic.
This paper may give better understanding for the proof transformation
and better resulting proofs of the proof transformation.

% (2)
In \cite{Ikebuchi25},
they simplified the proof transformation given in \cite{Berardi17}
by using ranking function synthesis given in \cite{Lee09}.
Their proof technique is that
they represent path relations by size-change graphs,
and then
they reduce
the induction principle for path relations
to
the induction principle for lexicographic ordering on natural numbers.
They rely on ranking functions provided in \cite{Lee09}.
In partiular, 
they used both Ramsey Theorem and Buchi automata in the meta-theory.

In this paper,
we use a similar proof technique to \cite{Ikebuchi25}.
In the same way as \cite{Ikebuchi25},
we will reduce
the induction principle for path relations
to
the induction principle for lexicographic ordering on natural numbers.
However, instead of depending on 
ranking function synthesis provided in \cite{Lee09},
we will give intuitionistic 
ranking function synthesis provided in this paper.
Namely, we will show
the existence of a ranking function for size-change termination condition
in the fan-in-free case by using intuitionistic logic,
where a size-change graph is called fan-in free \cite{Benamram09} if
the following holds:
if any two arguments in a calling function have an edge to the same
argument of a called function, then the two arguments are the same.”

In this paper,
for completing intuitionistic proof transformation,
we will also present intuitionistic ranking function synthesis
for fan-in-free case.
First, we will use a sequent with sequences of formulas,
which are used for cyclic proof system $\CLKIDomega$ given in 
\cite{Brotherston06},
instead of a sequent with sets of formulas,
which are used for cyclic proof system $\CLKIDomega$ given in 
\cite{Brotherston11}.
Then a set of size-change graphs 
which describes the path relation becomes fan-in free.

In \cite{Benamram09,Benamram10},
they gave ranking function synthesis for fan-in-free case.
Their proofs are classical.
In this paper, we will present a simple and intuitionistic proof
for ranking function synthesis for fan-in-free case.

When we obtain intuitionistic ranking function synthesis,
it is straightforward to obtain intuitionistic proof transformation
from cyclic proofs to inductive proofs, since
the correctness proof for the proof transformation in \cite{Berardi17} 
is intuitionistic except 
the induction principle for the path relation.

Our techniques for intuitionistic ranking function synthesis is as follows.
(1) We propose a novel notion of a domain-stable sequence.
For a given sequence $R_1,\ldots,R_n$ of relations,
we call the sequence {\em domain-stable} if
$\exists y(x R_i y)$ implies
$\exists y(x R_i \circ R_{i+1} y)$ for any $i \in [1,n-1]$.
(2) We presents some notion of weight for any domain-stable sequence.
The use of the size-change termination condition in a whole proof is only for 
guaranteeing the well-definedness of the weight.
The weight decreases when the computation goes with a domain-stable sequence.
(3) We provide an algorithm to obtain a domain-stable sequence for
any (possibly non-domain-stable) given sequence. 
By this algorithm,
the weight decreases when the computation goes with any sequence.

Section 2 gives background for size-change graphs and ranking functions.
Section 3 defines a domain-stable sequence and
presents intuitionistic synthesis of a ranking functions
for fan-in-free case.
Section 4 give background for inductive proofs and cyclic proofs.
Section 5 presents intuitionistic proof transformation from
cyclic proofs to inductive proofs.
Section 6 concludes.